\documentclass[11pt,a4paper]{article}
\usepackage[margin=2.5cm]{geometry}
\usepackage{booktabs}
\usepackage{graphicx}
\usepackage{xcolor}
\usepackage{hyperref}
\usepackage{enumitem}
\usepackage{titlesec}
\usepackage{fancyhdr}
\usepackage{microtype}
\usepackage{parskip}
\usepackage{fontspec}
\usepackage{polyglossia}
\usepackage{amsmath}
\usepackage{amssymb}

\setmainlanguage{english}
\setotherlanguage{hebrew}
\newfontfamily\hebrewfont{DejaVu Sans}[Script=Hebrew]

\hypersetup{colorlinks=true, linkcolor=blue, urlcolor=blue}

\pagestyle{fancy}
\fancyhf{}
\rhead{Knesset Protocol Journal Project}
\lhead{Step 6 --- Annotation}
\rfoot{Page \thepage}
\lfoot{Tomer Morad, Or Israeli, Yoel Marko --- July 2026}

\titleformat{\section}{\large\bfseries}{\thesection.}{0.6em}{}[\titlerule]
\titleformat{\subsection}{\normalsize\bfseries}{\thesubsection}{0.6em}{}

\begin{document}

\begin{center}
  {\LARGE \textbf{Annotation: Full-Protocol Gold Segmentation}}\\[0.3em]
  {\large Step 6 --- Milestone 1: the annotated evaluation set}\\[0.6em]
  {\normalsize Tomer Morad \quad Or Israeli \quad Yoel Marko}\\
  {\normalsize July 2026}
\end{center}

\vspace{0.5em}

\section{Objective}

\textbf{Milestone 1} requires a manually-labeled evaluation set before any quantitative
scoring is meaningful. The original plan (this step's own initial scope) was a 100-pair
sample of span pairs labeled \texttt{same}/\texttt{related}/\texttt{new}, built on Step 2's
marker-based topic spans. That plan was superseded mid-project (see \S3) by a more ambitious
target: manually re-segment every protocol into fine-grained topic spans, which yields both
Task~1 gold (segment boundaries) and Task~2 gold (recurring-topic chains) simultaneously and
at full corpus coverage, rather than a small sample.

\section{Method}

A browser-based annotation tool (\texttt{code/serve.py}, Flask) shows each protocol's
transcript with auto-detected \texttt{<{}<} \texthebrew{נושא} \texttt{>{}>} marker boundaries pre-filled, and lets
an annotator confirm, edit, split, or merge segments, labeling each with a topic string reused
across sessions where applicable.

Coverage started at 21/213 protocols (human-annotated by the three of us). The remaining
190/213 were completed by an LLM (Claude), reading each protocol's full utterance list to find
genuine subject changes rather than relying on the single auto-detected marker per protocol —
critical because many sessions (especially year-end ``budget surplus reallocation'' omnibus
meetings) cover a dozen unrelated ministry line-items under one nominal agenda title. Every
LLM-added segment is tagged \texttt{"annotator": "llm"} at the segment level, kept fully
distinguishable from the 21 human-annotated protocols.

A 14-protocol subset flagged as highest-risk (over 1000 utterances, yet only 1--2 segments
assigned) was independently re-verified with exhaustive, non-sampled re-reading. 13/14 held up
unchanged; one (\texttt{25\_ptv\_2623525}) had a missed opening segment, corrected.

\section{Results}

\begin{center}
\small
\begin{tabular}{lrp{6.5cm}}
\toprule
\textbf{Metric} & \textbf{Value} & \textbf{Note} \\
\midrule
Protocols segmented & 211 / 213 & 2 unsegmented (500K+ char marathon sessions, deferred) \\
Human-annotated protocols & 21 & Original team annotation \\
LLM-annotated protocols & 190 & Tagged \texttt{annotator: llm} per segment \\
Total segments & 522 & Across all 211 protocols \\
High-risk re-verification & 13/14 unchanged & 1 correction found and applied \\
\bottomrule
\end{tabular}
\end{center}

A separate 100-pair \texttt{same}/\texttt{related}/\texttt{new} set was also produced (LLM-judged,
content-only, blind to sampling stratum) as originally planned, before the full-segmentation
approach superseded it. It is preserved as \texttt{outputs/annotations\_llm.json} /
\texttt{gold\_labels.json} but is \textbf{not} used by any downstream evaluation script --- see
Step 7, which is built directly on the full segmentation instead.

\section{What Needs Improvement}

\begin{enumerate}[leftmargin=1.4em]
  \item \textbf{2 protocols remain unsegmented} --- both marathon overnight budget sessions
        (783K and 574K characters) using a legacy \texttt{<{}<} \texthebrew{הצח} \texttt{>{}>}-only marker format
        that the standard parser mishandles. Needs a dedicated chunked-reading pass.
  \item \textbf{No label canonicalization.} Near-duplicate topic strings across protocols are
        not merged (e.g. two segments about the same matter with slightly different generated
        titles will not be detected as recurring). Accepted per the project's own asymmetric-cost
        rule (false-split $\ll$ false-merge), but it means Step 7's ``22 recurring topics'' is
        very likely an \emph{undercount} of true recurrence.
  \item \textbf{The 100-pair \texttt{same}/\texttt{related}/\texttt{new} set is orphaned.} It
        has richer 3-way labels than the full-segmentation approach's binary same/not-same
        chains, and could still be mined for calibration examples (e.g.\ Step 8's few-shot
        prompts), but is not currently used anywhere in the pipeline.
  \item \textbf{Only spot-checked, not exhaustively verified.} 190 LLM-annotated protocols
        were produced; only the highest-risk 14 were re-verified with exhaustive reading. The
        remaining 176 rely on the original (partly sampled) reading pass.
\end{enumerate}

\vspace{1.5em}
\hrule
\vspace{0.4em}
{\small
\textbf{Artifacts (this step):} \newline
\texttt{outputs/segmentation.json} (211/213 protocols, 522 segments) \newline
\texttt{outputs/annotations\_llm.json}, \texttt{outputs/gold\_labels.json} \newline
\textbf{Code:} \texttt{code/serve.py}, \texttt{code/annotate.py}, \texttt{code/adjudicate.py},
\texttt{code/sample\_pairs.py}.
}

\end{document}
